\subsection{Вынужденные колебания: резонансный случай.}
Будем говорить что система входит в резонанс на $w^*$, если $|W_{jk}(w^*)| \ra +\infty$.
Это происходит если $\det D = 0$, то есть $\det(-A w^2 + iB w + C) = 0$.
Эквивалентно, среди корней $\det(A\lambda^2 + B\lambda + C)$ есть чисто мнимые. \\
Соответственно, если положение равновесия системы асимптотически устойчиво, то в этой системе не может быть резонансных случаев. \\
Но, если, всё-таки случай резонансный, то решение из прошлого параграфа непригодно и частное решение нужно искать в другом виде.
Для линейной консервативной системы мы можем записать уравнения в виде
\[
    A\ddot{\bt{q}} + C\bt{q} = \bt{Q}(t).
\]
После перехода к нормальным координатам $\bt{q} = U\bt{\theta}$ мы получим
\[
    \ddot{\bt{\theta}} + \Lambda \bt{\theta} = U^{T}\bt{Q} = U^T \bt{a}\cos wt.
\]
В частности, для каждого из осцилляторов
\[
    \ddot{\theta_k} + \omega^2_k \theta_k = u_k^T \bt{a} \cos wt.
\]
Если $w \neq w_k$, то
\[
    \theta_k(t) = \dfrac{\bt{u}_k^T \bt{a}}{w^2_k - w^2}\cos wt.
\]
Если $w = w_s$:
\[
    \theta_s(t) = \dfrac{\bt{u}_s^T \bt{a}}{2w_s}t \sin w_s t.
\]
(Если резонанс 1-кратности).
\begin{note}
    Следует заметить, что совпадение частоты вынуждающей силы $w$ с одной из собственных частот $w_s$ не вызовет резонанса если $\bt{u}_s^T\bt{a} = 0$.
\end{note}
\begin{note}
    Следует заметить, что в системе будет резонанс на единственной частоте $w_k$ если $\bt{a} = \lambda A \bt{u}_k$.
\end{note}
\subsection{Параметрический резонанс.}

\begin{definition}
    Параметрическим резонансом лагранжевой системы будем называть явление потери устойчивости положения равновесия при наложении периодического по времени возмущения.
\end{definition}
Существует несколько классических примеров систем с параметрическим резонансом:
\begin{itemize}
    \item Уравнение Хилла: $\ddot{x} + w^2(t) x = 0$, где $w(t)$ -- периодическая функция с некоторым периодом $T$.
    \item Уравнение Матье (частный случай уравнения Хилла): $\ddot{x} + w_0^2(1 + h \cos \gamma t)x = 0$.
    \item Осциллятор Ван-дер-Поля: $\ddot{x} - \mu(1 - x^2)\dot{x} + x = 0$.
    \item Осциллятор Дюффинга: $\ddot{x} + \delta \dot{x} + \alpha x + \beta x^3 = \gamma \sin w t$.
\end{itemize}
Рассмотрим поподробнее уравнение Хилла.
Пусть $X(t)$ -- фундаментальная матрица системы.
Тогда, поскольку матрица системы имеет вид
\[
    A(t) =
    \begin{pmatrix*}
        0 & 1 \\
        -w^2(t) & 0
    \end{pmatrix*}
\]
То по формуле Лиувилля-Остроградского
\[
    \dot{W} = \tr A(t) W = 0.
\]
И вронскиан системы сохраняется.
В то же время, $X(0) = I$, а следовательно $\det X(t) = 1$.
\begin{definition}
    Матрицей монодромии будет называть фундаментальную матрицу системы, вычисленную за период.
\end{definition}
Заметим, что в таком случае значение фундаментальной матрицы в $t + T$ можно вычислить при помощи матрицы монодромии:
\[
    X(t + T) = X(t)M.
\]
Это действительно так, поскольку если $x_1(t)$ и $x_2(t)$ -- решения системы, то и $x_1(t + T), x_2(t + T)$ -- тоже являются решениями системы.
Поэтому если мы рассмотрим $Y(t) = X(t + T)$, то
\[
    \dot{Y}(t) = \dot{X}(t + T) = A(t + T)X(t + T) = A(t)X(t).
\]
Кроме того, $Y(0) = X(T) = M$.
В тоже время, $X(t)$ -- решение задачи Коши (для матричного уравнения $\dot{X} = A(t)X$) с начальным условием $X(0) = I$.
Тогда решением поставленной задачи Коши будет $Y(t) = X(t)X^{-1}(0)Y(0) = X(t)M$.
В силу единственности решения и того, что $Y(t) = X(t + T)$ мы и получаем искомое соотношение:
\[
    X(t + T) = X(t)M.
\]
Из этого следует то, что поведение решений системы  $t \ra \infty$ определяется спектром матрицы монодромии.
Найдём собственные значения матрицы монодромии $M = X(T)$:
\[
    \det(M - \mu I) = \mu^2 - \tr M \mu + 1.
\]
Если $|\tr M| > 2$, то собственные значения $\mu_1$ и $\mu_2$ -- вещественны, и (без ограничения общности) $|\mu_1| > 1$ а $|\mu_2| < 1$.
А следовательно решение уравнения Хилла будет неустойчиво и возникнет <<параметрический резонанс>>. \\
Если же $|\tr M| < 2$, то собственные значения $\mu_1$ и $\mu_2$ будут комплексными с единичным модулем и решения будут устойчивыми.

\section{Классическая гамильтонова механика.}
\subsection{Преобразование Лежандра.}
\begin{definition}
    Векторное поле $\phi: U \to \R^n$ (где $U \subset \R^n$ -- область в $\R^n$) будем называть потенциальным, если существует достаточно гладкая (дважды непрерывной дифференцируемости достаточно) функция $V$ такая, что
    \[
        \phi = \dfrac{\partial V}{\partial\bt{x}}.
    \]
\end{definition}
\begin{definition}
    Преобразованием Лежандра будем называть невырожденное потенциальное векторное поле, то есть $\phi$ такое, что
    \[
        \det D\phi \neq 0.
    \]
    (Здесь и далее $D\phi$ - матрица Якоби отображения $\phi$).
\end{definition}
\begin{note}
    Заметим, что требование невырожденности векторного поля в точности совпадает с требованием невырожденности гессиана его потенциала:
    \[
        \det \Hess V \neq 0.
    \]
\end{note}
\begin{theorem}\label{th:1.1}
    Пусть $\bt{y} = \phi(\bt{x})$~---~преобразование Лежандра с потенциалом $V$.
    Тогда в окрестности каждой точки $\bt{x}^0$ существует обратная замена координат $\bt{x} = \psi(\bt{y})$ которая также является преобразованием Лежандра.
    Новый потенциал $W$ получается как
    \[
        W(\bt{y}) = \big(\bt{x} \cdot \bt{y} - V(\bt{x})\bigr|_{\bt{x} = \psi(y)}.
    \]
\end{theorem}
\begin{proof}
    Заметим, что в силу невырожденности векторного поля $\phi$ оно является локально обратимым и в окрестности каждой точки $\bt{x}^0$ определена обратная к ней $\psi(\bt{y})$ (по теореме о локальном диффеоморфизме), причём якобиан $\psi$ также не равен 0. \\
    Покажем потенциальность $\psi$.
    Для этого рассмотрим функцию \[W(\bt{y}) = \big( \bt{x} \cdot \bt{y} - V(\bt{x}) \bigr|_{\bt{x} = \psi(\bt{y})}.\]
    Покажем что она действительно является потенциалом для $\psi$:
    \[
        \dfrac{\partial W}{\partial \bt{y}} = \dfrac{\partial \psi}{\partial \bt{y}} \cdot \bt{y} + \psi(\bt{y}) - \dfrac{\partial V}{\partial \psi(\bt{y})} \cdot \dfrac{\partial \psi}{\partial \bt{y}} = \dfrac{\partial \psi}{\partial \bt{y}} \cdot \bt{y} + \psi(\bt{y}) - \bt{y} \cdot \dfrac{\partial \psi}{\partial \bt{y}} = \psi(\bt{y}).
    \]
    Доказательство теоремы завершено.
\end{proof}
\begin{note}
    Нетрудно заметить инволютивность преобразования Лежандра, то есть после второго его применения мы вернёмся к изначальной системе координат.
\end{note}
\subsection{Канонические уравнения Гамильтона.}
Рассмотрим некоторую $n$-мерную лагранжеву динамическую систему с лагранжианом $L(\bt{q}, \dot{\bt{q}}, t)$.
Уравнения движения этой механической системы подчиняются уравнениям Лагранжа:
\[
    \biggr(\dfrac{d}{dt}\dfrac{\partial}{\partial \dot{\bt{q}}} - \dfrac{\partial}{\partial \bt{q}}\biggr)L = 0.
\]
Пусть $\bt{p} = \frac{\partial L}{\partial \dot{\bt{q}}}$.
Здесь и далее будем называть $\bt{p}$ обобщёнными импульсами. \\
Рассмотрим замену координат
\[
    \phi: (\bt{q}, \dot{\bt{q}}, t) \mapsto (\bt{q}, \bt{p}, t).
\]
Она является преобразованием Лежандра и потенциалом по \hyperref[th:1.1]{теореме} будет $H$, называемый гамильтонианом, и равный
\[
    H = \big\langle \bt{p} \cdot \dot{\bt{q}} - L(\bt{q}, \dot{\bt{q}}, t) \big|_{\dot{\bt{q}} = \psi(\bt{q}, \bt{p}, t)}.
\]
где $\psi(\bt{q}, \bt{p}, t)$~---~обратная замена координат. \\
При такой замене координат уравнения Лагранжа переходят в канонические уравнения Гамильтона:
\[
        \dot{\bt{q}} = \dfrac{\partial H}{\partial \bt{p}}, \quad
        \dot{\bt{p}} = -\dfrac{\partial H}{\partial \bt{q}}.
\]
\begin{note}
    Стоит заметить, что мы можем использовать преобразование Лежандра только если $\det \dfrac{\partial^2 L}{\partial \dot{\bt{q}} \partial \dot{\bt{q}}} \neq 0$ -- то есть если векторное поле $\phi$, порождённое лагранжианом, невырождено.
\end{note}
\begin{note}
    Для консервативных систем $H$ совпадает с полной энергией системы (где обобщённые скорости выражены через обобщённые импульсы).
\end{note}
\begin{example}
    Пускай у нас есть одномерный осциллятор с лагранжианом
    \[
        L = \dfrac{m \dot{x}^2}{2} - \dfrac{mw^2x^2}{2}.
    \]
    Найдём обобщённый импульс:
    \[
        p = \dfrac{\partial L}{\partial \dot{x}} = m \dot{x}.
    \]
    (В этом частном случае обобщённый импульс совпадает с обычным -- что и является причиной подобного названия). \\
    Теперь запишем гамильтониан:
    \[
        H = p\dot{x} - L = \dfrac{p^2}{m} - (\dfrac{p^2}{2m} - \dfrac{mw^2 x^2}{2}) = \dfrac{p^2}{2m} + \dfrac{mw^2 x^2}{2}.
    \]
    Канонические уравнения Гамильтона будут иметь следующий вид:
    \begin{equation*}
            \dot{x} = \dfrac{\partial H}{\partial p} = \dfrac{p}{m}, \quad
            \dot{p} = -\dfrac{\partial H}{\partial x} = - mw^2 x.
    \end{equation*}
\end{example}
\begin{note}[Прим. техающего; Не помню, было ли это на лекции]
    Обычно функция Лагранжа неавтономной механической системы имеет вид $L = L_2 + L_1 + L_0$,
    где \[L_0 = -V(\bt{q}, t), \quad L_1 = \langle \bt{a}(\bt{q}, t), \dot{\bt{q}} \rangle, \quad L_2 = \dfrac{1}{2}\langle A(\bt{q}, t)\dot{\bt{q}}, \dot{\bt{q}} \rangle\]
    Тогда обобщённые импульсы выражаются как
    \[
        \dot{p} = A(\bt{q}, t)\dot{\bt{q}} + a(\bt{q}, t), \quad \dot{\bt{q}} = A^{-1}(\bt{q}, t)(\bt{p} - \bt{a}(\bt{q}, t)).
    \]
    И, соответственно, функция Гамильтона будет иметь вид
    \[
        H(\bt{q},\bt{p}, t) = \dfrac{1}{2}\langle A^{-1}(\bt{q}, t)(\bt{p} - \bt{a}(\bt{q}, t)), \bt{p} - \bt{a}(\bt{q}, t) \rangle + V(\bt{q}, t).
    \]
\end{note}
\begin{proposition}
    Если $q_k$ (или $p_k$) -- циклические координаты, то сопряжённая им координата сохраняется. \\
    Если циклической координатой является время, то сохраняется сам гамильтониан.
\end{proposition}
\begin{proof}
    Первая часть утверждения тривиально следует из канонических уравнений Гамильтона:
    \[
        \dot{p_k} = -\dfrac{\partial H}{\partial q_k}, \quad \dot{q_k} = \dfrac{\partial H}{\partial p_k}.
    \]
    Тогда, если $q_k$ -- циклическая, то есть $\dfrac{\partial H}{\partial q_k} = 0$, то $\dot{p}_k = 0$ и $p_k = const$ на решения системы.
    Аналогичное верно и для $p_k$. \\
    Теперь рассмотрим производную $H$ по времени:
    \[
        \dfrac{dH}{dt} = \dfrac{\partial H}{\partial t} + \dfrac{\partial H}{\partial \bt{q}} \dot{\bt{q}} + \dfrac{\partial H}{\partial \bt{p}}\dot{\bt{p}} = \dfrac{\partial H}{\partial t} + \dfrac{\partial H}{\partial \bt{q}} \dfrac{\partial H}{\partial \bt{p}} - \dfrac{\partial H}{\partial \bt{p}} \dfrac{\partial H}{\partial \bt{q}} = \dfrac{\partial H}{\partial t}.
    \]
    Отсюда видно, что если $\dfrac{\partial H}{\partial t} = 0$, то и $\dot{H} = 0$, а значит $H = const$ на решениях системы.
\end{proof}